\section{Synchronized Reference Trajectory Generation}
\label{sec:initialization}
A coordinated reference determines the initial corridors and supplies a warm start for optimization. Static A* guidance~\cite{hart1968} provides routes through the environment, and a synchronous social force model (SFM)~\cite{helbing1995} establishes an initial passing pattern. Parameterized social-force references can guide the selection of cooperative maneuvers~\cite{li2023priority}. The same scalar timing rule places every geometric path on the prescribed rest-to-rest horizon.

\subsection{Geometric Guidance and Cooperative Interaction}
For each UAV, a collision-free start--goal segment is used when available. Otherwise, a three-dimensional A* search in the inflated obstacle map generates a guide path, which is shortened by line-of-sight checks. The inflation includes the UAV footprint, clearance, and a margin for inter-node motion.

Let $\mathbf q_i$ be a point a look-ahead distance $\ell_g$ along the guide, and let $d_i=\norm{\mathbf g_i-\bp_i}_2$ be the distance to the goal. A preferred speed $v_{\rm ref}$, goal-approach time scale $\tau_g$, and velocity-response time scale $\tau_v$ define
\begin{equation}
 \bv_i^{\rm pref}=\min\{v_{\rm ref},d_i/\tau_g\}
 \frac{\mathbf q_i-\bp_i}{\max\{\norm{\mathbf q_i-\bp_i}_2,\ell_g\}},
 \qquad \mathbf f_i^{\rm att}=\frac{\bv_i^{\rm pref}-\bv_i}{\tau_v}.
 \label{eq:attraction}
\end{equation}
The guide point advances monotonically along the route, and the preferred speed decreases near the goal.

For neighboring UAVs, define $d_{ij}=\norm{\bp_i-\bp_j}_2$ and $\widehat{\mathbf r}_{ij}=(\bp_i-\bp_j)/d_{ij}$. For $\epsilon_d<d_{ij}<d_{\rm cut}$, the pairwise response is
\begin{equation}
 \mathbf f_{ij}=A_p\exp\!\left[
 \min\left\{\chi_{\max},\frac{2r+c+b_p-d_{ij}}{B_p}\right\}\right]
 \left(\widehat{\mathbf r}_{ij}+\kappa\,\mathbf e_z\times\widehat{\mathbf r}_{ij}\right).
 \label{eq:repulsion}
\end{equation}
Here $A_p$ sets the repulsion strength, $B_p$ its decay length, $b_p$ an interaction buffer, and $\chi_{\max}$ a numerical cap. Outside this distance range the pair response is zero. In particular, the small guard $\epsilon_d$ suppresses normalization at coincident positions. The coefficient $\kappa$ controls the tangential component, which encourages a consistent passing direction around the common vertical axis $\mathbf e_z$. The paired contributions are equal and opposite before goal tapering; their directions depend on geometry.

For compact notation, set $\psi(d;s)=\exp[\min\{\chi_{\max},(s-d)/B_p\}]$. The nearest point on obstacle $b$ is $\mathbf c_{ib}=\clip(\bp_i,\underline{\mathbf o}_b,\overline{\mathbf o}_b)$, with $d_{ib}=\norm{\bp_i-\mathbf c_{ib}}_2$. The obstacle and wall responses are
\begin{subequations}\label{eq:obstacle-forces}\begin{align}
 \mathbf f_i^{\rm obs}
 &=A_o\!\sum_{b:\,\epsilon_d<d_{ib}<d_o}
 \psi(d_{ib};s_o)\frac{\bp_i-\mathbf c_{ib}}{d_{ib}},
 \qquad s_o=r+c+v_{\max}h/2+\epsilon_g,\\
 f_{i,d}^{\rm wall}
 &=A_o\big[\psi(\ell_{i,d}^{-};b_w)-\psi(\ell_{i,d}^{+};b_w)\big],\\
 \ell_{i,d}^{-}&=p_{i,d}-\underline w_d-r,\qquad
 \ell_{i,d}^{+}=\overline w_d-r-p_{i,d}.
\end{align}\end{subequations}
Here $A_o$ is the response gain, $d_o$ the obstacle interaction range, and $b_w$ the boundary-response buffer. Obstacle directions within the zero-distance guard are omitted. To permit arrival at a separated goal, the summed repulsion is weighted by $\theta_i=\min\{1,d_i/d_g\}^2$, where $d_g$ specifies the goal-approach region. The total acceleration is
\[
 \mathbf f_i=\mathbf f_i^{\rm att}
 +\theta_i\left(\sum_{j\ne i}\mathbf f_{ij}
 +\mathbf f_i^{\rm obs}+\mathbf f_i^{\rm wall}\right).
\]
Coordinate caps $a_{\max}/\sqrt3$ and $v_{\max}/\sqrt3$ provide a conservative numerical step for the geometric simulation. Clipping this response to the acceleration caps and the next-step velocity caps gives $\mathbf f_i^{\rm clip}$. All UAVs advance synchronously:
\begin{equation}
 \bp_i^+=\bp_i+\Delta t\,\bv_i+\tfrac12\Delta t^2\mathbf f_i^{\rm clip},
 \qquad \bv_i^+=\bv_i+\Delta t\,\mathbf f_i^{\rm clip}.
 \label{eq:simulation}
\end{equation}
When $d_i<\epsilon_p$ and $\norm{\bv_i}_2<\epsilon_v$, the simulated state is set to $(\mathbf g_i,\mathbf0)$ and held there while the remaining UAVs proceed. This gives exact goal knots with a common sequence of time indices. If the cooperative simulation reaches its limit, the static guide paths provide a fallback geometric reference. The same fallback is available to every guided method. A missing guide or an inadmissible fallback time law terminates reference generation. The SFM supplies coordinated geometric references; collision clearance and motion feasibility are established by the subsequent optimization and acceptance stages.

\subsection{Common-Horizon Rest-to-Rest Timing}
For UAV $i$, remove repeated consecutive points from its simulated geometric polyline. Let $\mathbf z_i(q)$ be its piecewise-linear arc-length parameterization, with $q\in[0,L_i]$. Each nonzero segment has $\norm{\mathbf z_i'(q)}_2=1$. The scalar profile is constructed on the optimization grid:
\begin{subequations}\label{eq:scalarprofile}\begin{align}
 b_k&=\min\{v_{\max},a_{\max}t_k,a_{\max}(T_f-t_k)\},\qquad
 Z=\frac h2\sum_{k=0}^{K-1}(b_k+b_{k+1}),\\
 u_{i,k}&=\frac{L_i}{Z}b_k,\qquad
 q_{i,k}^{\star}=\frac h2\sum_{\ell=0}^{k-1}(u_{i,\ell}+u_{i,\ell+1}),\\
 \alpha_{i,k}&=(u_{i,k+1}-u_{i,k})/h .
\end{align}\end{subequations}
The profile satisfies zero endpoint speeds and covers exactly $L_i$ by $T_f$. Its maximum speed and absolute acceleration must satisfy $v_{\max}$ and $a_{\max}$. An SFM route that exceeds this capacity triggers the guide fallback; a fallback or straight route that still exceeds it ends the attempt. The horizon is unchanged.
The reference states are
\begin{equation}
 \widehat{\bp}_{i,k}=\mathbf z_i(q_{i,k}^{\star}),\qquad
 \widehat{\bv}_{i,k}=u_{i,k}\mathbf z_i'(q_{i,k}^{\star}).
 \label{eq:reference}
\end{equation}
At a geometric knot, the outgoing tangent is used, with the incoming tangent at the terminal knot. The common timing rule gives each variant the same horizon and boundary-speed conditions. Route-specific arc-length scaling retains the simulated paths but can change their relative passage times; cooperative optimization resolves the remaining temporal conflicts. Bends can also produce vector-dynamics defects. The scalar law establishes reference timing, while final motion feasibility is determined after optimization and reconstruction.
